\documentclass[11pt,a4paper]{article}
\usepackage[utf8]{inputenc}
\usepackage[T1]{fontenc}
\IfFileExists{polish.ldf}{\usepackage[polish]{babel}}{}
\IfFileExists{lmodern.sty}{\usepackage{lmodern}}{}
\usepackage[margin=26mm]{geometry}
\usepackage[authoryear,round]{natbib}
\usepackage{amsmath}
\usepackage[hidelinks]{hyperref}
\title{Struktura czynnikowa i trafność narzędzia DERA-HRI do oceny doświadczenia interakcji z humanoidem}
\author{Mateusz Pomianek}
\date{MSC-DERA-P01 | Protokół prospektywny 1.0.0 | 27 listopada 2026}
\begin{document}
\maketitle
\begin{abstract}
Badanie ma dostarczyć dowodów trafności interpretacji wyników DERA-HRI po bezpośredniej interakcji z humanoidem. Projekt łączy ocenę treści i procesu odpowiadania, eksplorację na pierwszej próbie oraz konfirmację na niezależnej próbie. Wyniki będą odnoszone do wersji kwestionariusza, języka i zdefiniowanego scenariusza. Przed rozpoczęciem badań niezbędne jest ustalenie proweniencji pozycji, klucza oraz planu użycia narzędzia.
\end{abstract}
\section{Wprowadzenie}
Kwestionariusz doświadczenia po interakcji może operacjonalizować reakcję na konkretną sytuację. Miary Godspeed obejmują między innymi postrzegane bezpieczeństwo i sympatię, a RoSAS odnosi się do ocen ciepła, kompetencji i dyskomfortu \citep{bartneck,carpinella}. Plan walidacji dobiera komparatory do rzeczywistej treści podskal. Ciepło i sympatia mogą współwystępować z małym dyskomfortem, lecz są odrębnymi hipotezami pomiarowymi. Klasyfikacja kontekstu HRI pozwala opisać rolę człowieka i robota, zadanie oraz ich zależność \citep{onnasch}. Projekt MSC-HRI-P08 bada relację między subiektywnym a fizycznym bezpieczeństwem. P01 dotyczy znaczenia i jakości wyników kwestionariusza w zdefiniowanym użyciu.

\section{Pytania badawcze}
Czy pozycje pokrywają wymagany obszar doświadczenia i czy respondenci odnoszą odpowiedzi do zakończonej interakcji? Czy model wybrany w fazie eksploracyjnej odtwarza zależności między wskaźnikami w niezależnej próbie? Z jaką precyzją można oszacować rzetelność wyników oraz zależności z dopasowanymi i odmiennymi miarami? Czy progi i ładunki są porównywalne w scenariuszach, dla których planowane są porównania grup?

Po ustaleniu klucza prerejestracja określi liczbę i znaczenie czynników. Hipoteza trafności przewiduje silniejsze zależności latentne dla komparatorów zgodnych treściowo niż dla odmiennych. Jeżeli wersja narzędzia obejmuje dyskomfort, jego wynik powinien dodatnio korelować z dopasowaną domeną RoSAS i ujemnie z postrzeganym bezpieczeństwem Godspeed. Wyniki urządzenia i deklaracje uczestników wymagają różnych źródeł dowodów.

\section{Materiał i metody}
Faza A obejmuje kontrolę praw do użycia pozycji, ocenę ekspertów oraz wywiady poznawcze z dorosłymi polskojęzycznymi. Bank dalszych pozycji i reguły punktowania wynikają z potwierdzonej wersji. Faza B obejmuje eksplorację po porównywalnej interakcji zadaniowej z humanoidem. Faza C przewiduje odrębną rekrutację, blokadę modelu, pomiar kotwic i analizę potwierdzającą. Kolejność skal po interakcji randomizuje się. Przerwania interakcji, różnice między operatorami i odstępstwa sprzętowe rejestruje się jako informacje o realizacji procedury. Analiza pierwszorzędowa wykorzystuje pierwszą udokumentowaną sesję każdego uczestnika.

\section{Plan analizy}
Dla pozycji $j$ z odpowiedzią porządkową i progami $\tau_{jk}$ model zakłada
\[
 \Pr(Y_{ij}\leq k\mid\boldsymbol\eta_i)
 =\Phi(\tau_{jk}-\boldsymbol\lambda_j^\top\boldsymbol\eta_i).
\]
Pierwszorzędowy wskaźnik: replikacja zamrożonego modelu progowo-czynnikowego w niezależnej próbie, oceniana łącznie przez dopasowanie globalne, lokalne reszty i stabilność ładunków. Eksploracja wykorzystuje korelacje polichoryczne, analizę równoległą i rotację ukośną. Dobór estymatora uwzględnia porządkowy charakter danych \citep{flora}. Na etapie konfirmacji ocenia się wszystkie prerejestrowane modele konkurencyjne, w tym model jednoczynnikowy i model skorelowanych czynników. Modele wyższego rzędu wymagają uzasadnienia treściowego i identyfikowalności. Raport obejmuje zbieżność, rozwiązania niewłaściwe, rozkłady progów, ładunki oraz dopasowanie globalne i lokalne.

Omega oparta na modelu dotyczy każdej uzasadnionej podskali. Omega dla latentnej odpowiedzi i rzetelność obserwowanej sumy ordinalnej wymagają osobnego nazwania. W modelu jednoczynnikowym bez skorelowanych błędów i przy wariancji czynnika równej jeden stosunek dla sumy odpowiedzi latentnych ma postać $(\sum_j\lambda_j)^2/[(\sum_j\lambda_j)^2+\sum_j\theta_j]$. Niepewność estymuje się przez powtórne dopasowania bootstrapowe na poziomie osób. Inwariancja progów i ładunków wymaga jawnych ograniczeń identyfikujących \citep{wu}. Porównania średnich latentnych są warunkowe względem adekwatności modelu pomiarowego. Zamrożenie estymandu, modeli, kryteriów i reguł analizy nastąpi przed odczytem próby konfirmacyjnej.

\section{Wyniki}
Nie przeprowadzono rekrutacji ani pomiarów. Do planowania wykonano syntetyczną symulację ryzyka rzadkich kategorii dla hipotetycznego banku. W 24\,000 replikacji dla hipotetycznych 12 pozycji oceniono częstość wystąpienia co najmniej jednej kategorii z mniej niż pięcioma odpowiedziami. Przy próbie $N=400$ i prawdopodobieństwie najrzadszej kategorii 0,02 odsetek wyniósł 66,6\%, a przy $N=600$ spadł do 8,15\%. Przy prawdopodobieństwie 0,01 pozostał na poziomie 97,0\% dla $N=600$. Pełną analizę liczebności dla replikacji modelu przeprowadza się po rozpoznaniu struktury i progów autoryzowanej wersji.

\section{Dyskusja i ograniczenia}
Zakres wniosków wyznaczają wersja kwestionariusza, kontekst, populacja i sposób realizacji interakcji. Stabilność testu po ponownej ekspozycji zależy również od habituacji i zmiany doświadczenia. Interpretacja subiektywnego poczucia bezpieczeństwa wymaga rozróżnienia od wyników fizycznych. Kandydatem na czasopismo po walidacji jest \emph{International Journal of Social Robotics}.
\bibliographystyle{apalike}
\bibliography{references}
\end{document}
